\documentclass[11pt,a4paper]{article}
% ============================================================
% Shared preamble for the MCA104 CNIA exam-preparation notes.
%
% libertinus and tcolorbox are installed in the user TeX tree
% (/home/elish4h/texmf) via:
%   tlmgr --usermode install --ignore-warning \
%     tcolorbox libertinus libertinus-type1 libertinus-otf \
%     libertinus-fonts libertinust1math
%
% The box styling matches the repository's notes-template.tex.
% ============================================================

\usepackage[margin=0.72in]{geometry}
\usepackage{libertinus}
\usepackage{microtype}
\usepackage[T1]{fontenc}
\usepackage{xcolor}
\usepackage{enumitem}
\usepackage{tabularx}
\usepackage{array}
\usepackage{booktabs}
\usepackage{longtable}
\usepackage{amsmath}
\usepackage{amssymb}
\usepackage{fancyhdr}
\usepackage{titlesec}
\usepackage[skins,breakable]{tcolorbox}
\usepackage{hyperref}
\usepackage{multicol}
\usepackage{parskip}

\definecolor{navy}{HTML}{183153}
\definecolor{blue}{HTML}{2563EB}
\definecolor{green}{HTML}{166534}
\definecolor{red}{HTML}{B91C1C}
\definecolor{gold}{HTML}{92400E}
\definecolor{lightblue}{HTML}{EFF6FF}
\definecolor{lightgreen}{HTML}{F0FDF4}
\definecolor{lightred}{HTML}{FEF2F2}
\definecolor{lightgold}{HTML}{FFFBEB}
\definecolor{grey}{HTML}{475569}

\hypersetup{colorlinks=true,linkcolor=blue,urlcolor=blue}

\titleformat{\section}{\Large\bfseries\color{navy}}{}{0pt}{}
\titleformat{\subsection}{\large\bfseries\color{navy}}{}{0pt}{}
\titleformat{\subsubsection}{\normalsize\bfseries\color{blue}}{}{0pt}{}

\setlist[itemize]{leftmargin=*,nosep}
\setlist[enumerate]{leftmargin=*,nosep}

% Allow tiny line-length overruns to be absorbed rather than overfull.
\setlength{\emergencystretch}{2em}

\tcbset{
  parbox=false,
  enhanced,
  breakable,
  colbacktitle=blue,
  coltitle=white,
  fonttitle=\bfseries,
  boxrule=0.5pt,
  arc=1.5pt,
  left=6pt,right=6pt,top=4pt,bottom=4pt,
  before skip=6pt,after skip=6pt,
}

\newtcolorbox{exambox}{colback=lightblue,colframe=blue,title=Exam Focus}
\newtcolorbox{rememberbox}{colback=lightgold,colframe=gold,title=Remember}
\newtcolorbox{dangerbox}{colback=lightred,colframe=red,title=Do Not Confuse}
\newtcolorbox{answerbox}{colback=lightgreen,colframe=green,title=Answer Skeleton}
\newtcolorbox{pyqbox}{colback=white,colframe=grey,title=Past-paper question}
\newtcolorbox{solbox}{colback=lightgreen,colframe=green,title=Solution}

\newcommand{\mkmarks}[1]{\hfill\textbf{[#1 marks]}}
\newcommand{\cmd}[1]{\texttt{#1}}
\newcommand{\F}[1]{\ensuremath{#1}}
\newcommand{\subject}[1]{\lhead{#1}\rhead{MCA104 Exam Notes}\cfoot{\thepage}}

\subject{MCA104 -- Exam Overview and Past-Paper Analysis}

\title{\textbf{Computer Networks and Internet Architecture}\\}
\author{\large MCA104 Exam Preparation -- Overview, Past-Paper Analysis and Study Plan}
\date{Exam: Wednesday 18 November 2026}

\begin{document}
\maketitle

\begin{exambox}
\textbf{How to use this set.} Five module PDFs plus this overview. Every past-paper question we have is placed in the module it belongs to, so as you revise a module you can see exactly how it has been asked before. Section 6 is a five-day plan; Section 7 is the formula sheet.
\end{exambox}

\tableofcontents
\newpage

\section{The six papers we have}

\begin{center}\small
\begin{tabularx}{\textwidth}{>{\bfseries}p{0.26\textwidth} p{0.10\textwidth} X}
\toprule Paper & Marks & Format\\
\midrule
Sept 2026 CAE-I & 60 & 2 hours. Sections A, B, C, 10 marks each.\\
Oct 2026 CAE-II & 30 & 1 hour. Section A, B, C, 10 marks each.\\
Oct 2025 CAE-II & 20 & 60 minutes. Section A and B, 10 marks each.\\
Oct 2024 CAE-II & 20 & 60 minutes. Section A and B, 10 marks each.\\
Oct 2025 full & 75 & 3 hours. Section A: 5 questions $\times$ 15 marks, all compulsory.\\
Sept 2024 CAE-I & 40 & 90 minutes. Section A and B, 10 marks each.\\
\bottomrule
\end{tabularx}
\end{center}

\begin{exambox}
\textbf{The two most predictive papers} are the Sept 2026 CAE-I and the Oct 2026 CAE-II, because they are from your own batch and their coverage matches the current syllabus. Read Section 3 for those first.
\end{exambox}

\section{Marks by module}

\begin{center}\small
\begin{tabularx}{\textwidth}{>{\bfseries}p{0.14\textwidth} X p{0.16\textwidth}}
\toprule Module & Content & Papers hit\\
\midrule
1 & Data communication, OSI and TCP/IP, media, switching, multiplexing, error control, MAC & 6 of 6\\
2 & Network layer, addressing, subnetting, routing, IPv4, IPv6, ICMP & 6 of 6\\
3 & Transport layer, TCP, UDP, flow and congestion control, QoS & 6 of 6\\
4 & Application layer, DNS, DHCP, security, TLS, IPSec & 5 of 6\\
5 & Internet structure, ISP, IXP, AS, CDN, cloud & 5 of 6\\
\bottomrule
\end{tabularx}
\end{center}

\textbf{Module 1 carries the most marks overall} because both CAE-I papers are almost entirely Module 1. \textbf{Module 3 and Module 4 carry the most marks in the semester paper.}

\section{The recurring questions, module by module}

\subsection{Module 1}

\begin{center}\small
\begin{tabularx}{\textwidth}{>{\bfseries}p{0.16\textwidth} X}
\toprule Paper & Question\\
\midrule
Sept 2026 & Explain the importance of the layered approach.\\
Sept 2026 & Define data communication; five components; performance characteristics.\\
Sept 2026 & Random access protocols: ALOHA, slotted ALOHA, CSMA, CSMA/CD.\\
Sept 2026 & Transmission media: guided and unguided.\\
Sept 2026 & Switching techniques; compare circuit and packet.\\
Sept 2026 & TCP/IP suite versus OSI.\\
Sept 2026 & Multiplexing: FDM, TDM, WDM.\\
Sept 2026 & CRC for data 1101011011, generator 10011.\\
Oct 2025 & Define multiplexing; FDM and TDM with diagrams.\\
Oct 2025 & Data communication components; framing.\\
Oct 2025 & Go-Back-N versus Selective Repeat.\\
Oct 2025 & Flow control in the data link layer.\\
Oct 2025 & TDM versus FDM.\\
Oct 2025 CAE-II & CRC for 10011101 with $x^3+1$; CSMA/CD scenario with 5 computers.\\
Oct 2024 CAE-II & CRC for 10011101; sliding window protocols; LLC versus MAC.\\
Sept 2024 & Polar coding schemes; PCM; wireless versus wired; PSTN.\\
\bottomrule
\end{tabularx}
\end{center}

\subsection{Module 2}

\begin{center}\small
\begin{tabularx}{\textwidth}{>{\bfseries}p{0.16\textwidth} X}
\toprule Paper & Question\\
\midrule
Sept 2026 & RIP versus OSPF: algorithm, metric, convergence, scalability.\\
Sept 2026 & Structure of an IPv4 datagram with a diagram.\\
Oct 2025 & ICMP in troubleshooting: echo and destination unreachable.\\
Oct 2025 & Responsibilities of the network layer; IPv6 versus IPv4.\\
Oct 2025 & Functions of the network layer and its services to transport.\\
Oct 2025 & Routing tables and dynamic routing; subnetting.\\
Oct 2024 CAE-II & ARP for MAC mapping failure.\\
\bottomrule
\end{tabularx}
\end{center}

\subsection{Module 3}

\begin{center}\small
\begin{tabularx}{\textwidth}{>{\bfseries}p{0.16\textwidth} X}
\toprule Paper & Question\\
\midrule
Sept 2026 & Connection-oriented versus connectionless; TCP and UDP as examples.\\
Sept 2026 & Short notes: transport services, SSL/TLS, TCP, UDP.\\
Oct 2025 & Flow control mechanisms in TCP; ACKs and timers.\\
Oct 2025 & Connection-oriented versus connectionless; QoS parameters.\\
Oct 2025 & TCP versus UDP in real-time applications; why UDP wins.\\
Oct 2025 & How congestion occurs; how UDP keeps latency low.\\
Oct 2025 CAE-II & Go Back N, Selective Repeat, piggybacking.\\
\bottomrule
\end{tabularx}
\end{center}

\subsection{Module 4}

\begin{center}\small
\begin{tabularx}{\textwidth}{>{\bfseries}p{0.16\textwidth} X}
\toprule Paper & Question\\
\midrule
Oct 2026 CAE-II & HTTP, FTP and SMTP in supporting web, file transfer and email.\\
Oct 2026 CAE-II & Propose a security solution: cryptography, firewalls, VPNs.\\
Oct 2026 CAE-II & DNS and DHCP: name resolution and dynamic addressing.\\
Oct 2025 & Define cryptography; plaintext, ciphertext and key.\\
Oct 2025 & Application layer function; objective of SSL/TLS.\\
Oct 2025 & HTTP and FTP vulnerabilities and secure alternatives.\\
Oct 2025 & How firewalls filter packets; DNS resolution step by step.\\
Oct 2025 CAE-II & Protocol choice for web, email and file transfer; FTP active vs passive.\\
Sept 2024 & Short notes: HTTP, SMTP, UDP, ICMP.\\
\bottomrule
\end{tabularx}
\end{center}

\subsection{Module 5}

\begin{center}\small
\begin{tabularx}{\textwidth}{>{\bfseries}p{0.16\textwidth} X}
\toprule Paper & Question\\
\midrule
Oct 2026 CAE-II & Autonomous Systems and BGP.\\
Oct 2026 CAE-II & CDN for a far-away e-commerce site.\\
Oct 2026 CAE-II & Cloud computing: IaaS, PaaS, SaaS.\\
Oct 2025 & Cloud computing: role, scalability, cost.\\
Oct 2025 & Internet versus intranet; Tier-1, 2 and 3 ISPs.\\
Oct 2025 & Define an ISP and its types; redundancy; CDN definition.\\
\bottomrule
\end{tabularx}
\end{center}

\section{The one thing that is almost certain}

\textbf{CRC.} It appears in three of the six papers, and two of them use the identical numbers (data 10011101, generator $x^3+1$). Learn the division procedure cold; it is 10 easy marks.

\textbf{The TCP/UDP comparison.} Appears four times in different forms: as a table, as ``connection-oriented versus connectionless'', as ``why UDP for real time'', and as a short note. One table answers all four.

\textbf{FDM versus TDM.} Appears three times, always worth 3 marks.

\textbf{The TCP/IP versus OSI comparison.} Appears twice, always worth 10.

\section{Where the marks are}

\begin{center}\small
\begin{tabularx}{\textwidth}{>{\bfseries}p{0.26\textwidth} c X}
\toprule Type & Weight & Preparation\\
\midrule
Definition / list & 3 marks & The five components, the seven layers, the five QoS parameters, the five NIST characteristics.\\
Comparison table & 3--10 marks & OSI/TCP-IP, circuit/packet, FDM/TDM, TCP/UDP, RIP/OSPF, TLS/IPSec, IaaS/PaaS/SaaS.\\
Calculation & 10 marks & CRC is the only routine calculation. Subnetting and Dijkstra may appear.\\
Diagram & 10 marks & Layered models, encapsulation, TCP and IPv4 headers, three-way handshake, congestion sawtooth.\\
Scenario & 10 marks & The October 2026 style. Propose and justify; never just list products.\\
\bottomrule
\end{tabularx}
\end{center}

\section{Five-day plan, 13 to 17 November}

\begin{center}\small
\begin{tabularx}{\textwidth}{>{\bfseries}p{0.14\textwidth} p{0.10\textwidth} X}
\toprule Day & Module & Work\\
\midrule
Thu 13 Nov & 1 & Read Module 1 in full. Do both CRC questions by hand, twice, without looking. Draw the OSI and TCP/IP stacks from memory.\\
Fri 14 Nov & 2 & Read Module 2. Do the subnetting examples, then the Dijkstra example, then the fragmentation example.\\
Sat 15 Nov & 3 and 4 & Read Module 3, then Module 4. Write the TCP/UDP table from memory; then the TLS handshake in six steps.\\
Sun 16 Nov & 5 & Read Module 5. Write the CDN delivery process in five steps and the cloud service model table.\\
Mon 17 Nov & All & Only the last-minute page of each module, plus the formula sheet in Section 7. Do not read anything new.\\
\bottomrule
\end{tabularx}
\end{center}

\begin{exambox}
\textbf{Exam rules.} For a comparison question, draw the table first; prose loses marks. For the CRC, show every XOR step; the answer alone scores almost nothing. For a scenario question, name the threat, name the control, and say what the control prevents.
\end{exambox}

\section{Formula and fact sheet}

\subsection{Module 1}

\begin{center}\small
\begin{tabularx}{\textwidth}{>{\bfseries}p{0.26\textwidth} X}
\toprule Item & Formula or fact\\
\midrule
Nyquist capacity & $C = 2B\log_2 L$ bps\\
Shannon capacity & $C = B\log_2(1 + S/N)$\\
PCM bit rate & $R = f_s \cdot n$; voice $8000 \times 8 = 64$ kbps\\
Quantisation SNR & $\approx 6.02n + 1.76$ dB\\
GBN window & $N = 2^m - 1$\\
Selective Repeat window & $N = 2^{m-1}$\\
Sliding window utilisation & $U = N/(1+2a)$, $a = T_p/T_t$\\
Full utilisation & $N \ge 1 + 2a$\\
ALOHA throughput & $S = Ge^{-2G}$, max 18.4\%\\
Slotted ALOHA & $S = Ge^{-G}$, max 36.8\%\\
CSMA/CD condition & $T_t \ge 2T_p$\\
CSMA/CD efficiency & $U = T_t/(T_t + 2T_p)$\\
Full mesh links & $n(n-1)/2$\\
Hamming check bits & $2^r \ge m + r + 1$\\
\bottomrule
\end{tabularx}
\end{center}

\subsection{Module 2}

\begin{center}\small
\begin{tabularx}{\textwidth}{>{\bfseries}p{0.26\textwidth} X}
\toprule Item & Formula or fact\\
\midrule
Host bits & $h = 32 - n$\\
Addresses per subnet & $2^h$\\
Usable hosts & $2^h - 2$\\
Block size & $2^h$\\
Aggregation & $2^k$ contiguous prefixes $\Rightarrow$ subtract $k$\\
Fragment offset unit & 8 bytes\\
RIP metric & Hop count, max 15, 16 = unreachable\\
OSPF metric & Cost $= 10^8 / \text{bandwidth}$\\
Dijkstra & $O(n^2)$ or $O((n+e)\log n)$\\
Bellman-Ford & $D(v,X) = \min_w [c(v,w) + D(w,X)]$\\
\bottomrule
\end{tabularx}
\end{center}

\subsection{Modules 3 to 5}

\begin{center}\small
\begin{tabularx}{\textwidth}{>{\bfseries}p{0.26\textwidth} X}
\toprule Item & Formula or fact\\
\midrule
Well-known ports & 20/21 FTP, 22 SSH, 23 Telnet, 25 SMTP, 53 DNS, 67/68 DHCP, 80 HTTP, 110 POP3, 143 IMAP, 161 SNMP, 443 HTTPS\\
TCP header minimum & 20 bytes; options to 60\\
UDP header & 8 bytes\\
Effective send window & $\min(\text{rwnd}, \text{cwnd}) - $ bytes in flight\\
Slow start & Double \texttt{cwnd} per RTT\\
Congestion avoidance & $+1$ MSS per RTT\\
On timeout & \texttt{ssthresh} = \texttt{cwnd}/2, \texttt{cwnd} = 1 MSS\\
On 3 duplicate ACKs & Fast retransmit, halve \texttt{cwnd}\\
IntServ versus DiffServ & Per-flow reservation versus per-class marking\\
RSA key relation & $ed \equiv 1 \pmod{\phi(n)}$, $\phi = (p-1)(q-1)$\\
Signature & Hash, then encrypt with the \emph{private} key\\
TLS objectives & Confidentiality, integrity, authentication\\
IPSec modes & Transport (payload) and tunnel (whole packet)\\
DHCP & Discover, Offer, Request, Acknowledge\\
\bottomrule
\end{tabularx}
\end{center}

\section{The eight diagrams to be able to draw}

\begin{enumerate}
\item The seven-layer OSI stack and the four-layer TCP/IP stack side by side, with PDUs.
\item Encapsulation: message, segment, packet, frame, bits.
\item A frame format for CRC: data, appended zeros, remainder, transmitted codeword.
\item The TCP three-way handshake and the four-way teardown.
\item The IPv4 header as a row of labelled boxes.
\item The congestion sawtooth with slow start and congestion avoidance labelled.
\item The Internet hierarchy: access, Tier 3, Tier 2, Tier 1, IXP, content.
\item CDN delivery: DNS picks the nearest edge, cache hit or miss, origin.
\end{enumerate}

\section{Common mistakes that cost marks}

\begin{enumerate}
\item Using ordinary binary subtraction in a CRC division. It is XOR.
\item Forgetting to append $n$ zeros before dividing in CRC, where $n$ is the generator's degree.
\item Swapping the TCP/IP and OSI layer counts.
\item Saying ``FDM divides time''. FDM divides frequency; TDM divides time.
\item Describing an IXP as something that routes.
\item Saying Tier-1 ISPs buy transit. They do not.
\item Confusing flow control (\texttt{rwnd}, protects the receiver) with congestion control (\texttt{cwnd}, protects the network).
\item Writing a three-way teardown. TCP teardown is four-way.
\item Signing with the public key. Signing uses the private key.
\item Giving a comparison as two paragraphs instead of a table.
\item Answering a scenario question with a product list and no justification.
\item Forgetting that usable hosts $= 2^h - 2$, not $2^h$.
\end{enumerate}

\vfill
\begin{center}\textit{Every calculation in this set was computed and checked programmatically. The CRC remainders, the Dijkstra distances, the subnet boundaries and the RSA values are all verified.}\end{center}

\end{document}
